%% 绪论：算子与它的分解
%% 本文件由 tools/md2tex.py 自动生成，请勿直接编辑。

\chapter*{绪论\quad 算子与它的分解}
\addcontentsline{toc}{chapter}{绪论\quad 算子与它的分解}
\markboth{绪论}{绪论}


\section*{这一部分研究什么}
\addcontentsline{toc}{section}{这一部分研究什么}

算子是从一个向量空间映到它自身的线性映射：

\[
T:V\to V.
\]

因为输入与输出都在同一空间中，算子可以反复作用：

\[
v,\quad T(v),\quad T^2(v),\quad T^3(v),\ldots
\]

因此，算子不仅描述一次线性变化，也描述一个线性系统内部不断演化的结构。

本章研究的核心问题是：

\begin{jinsight}
能否将一个复杂算子拆成若干简单、彼此独立或具有明确耦合方式的部分？
\end{jinsight}

若可以找到一组适合的基，使算子的矩阵具有对角、块对角或 Jordan 块等简单形状，我们就能直接读出算子的长期作用、可逆性、特征值、幂的行为与不变结构。这对于物理系统的分析，微分方程的解，机器学习中权重矩阵的分析，研究线性方程组的解都是非常有用的。

\section*{为什么可以分解}
\addcontentsline{toc}{section}{为什么可以分解}

算子分解的共同基础是不变子空间。

若 $U$ 是 $V$ 的子空间，且：

\[
T(U)\subseteq U,
\]

则称 $U$ 是 $T$ 的不变子空间。

这意味着：一个向量只要从 $U$ 出发，无论施加 $T$ 多少次，都不会离开 $U$。

若存在：

\[
V=U_1\oplus\cdots\oplus U_m,
\]

并且每个 $U_j$ 都对 $T$ 不变，那么任意向量都能唯一写成：

\[
v=u_1+\cdots+u_m,
\qquad u_j\in U_j.
\]

由线性性：

\[
T(v)=T(u_1)+\cdots+T(u_m).
\]

每一部分仍留在自己的不变子空间中。因此，我们可以分别研究 $T$ 在每个 $U_j$ 上的作用，再将结果组合起来。

算子分解的本质，就是寻找足够多的不变子空间，并理解它们之间是否能够形成直和。

\section*{分解方式总览}
\addcontentsline{toc}{section}{分解方式总览}


\subsection*{一般实向量空间与复向量空间}
\addcontentsline{toc}{subsection}{一般实向量空间与复向量空间}

这一部分只使用线性结构，不要求内积。

\begin{jitem}
\item %
上三角化：将算子表示为沿一条不变子空间链逐层作用的形式；
\item %
对角化：将空间分解为一维不变子空间的直和；
\item %
Jordan 标准型：当特征向量不够时，用广义特征向量和 Jordan 链描述剩余耦合；
\item %
复化：将实算子扩展到复空间中，以暴露复特征值与复 Jordan 结构，再翻译回实空间中的二维不变块。
\end{jitem}

其中，对角化是 Jordan 标准型的特殊情形：每一个 Jordan 块都只有 $1\times1$ 时，Jordan 矩阵就是对角矩阵。

\subsection*{内积空间中的算子分解}
\addcontentsline{toc}{subsection}{内积空间中的算子分解}

内积为向量空间增加长度、正交与伴随结构。

对于自伴算子、正规算子、酉算子等与内积相容的特殊算子，谱定理能够给出规范正交的特征向量基。此时算子不仅可对角化，而且可以在保持长度和正交关系的坐标中对角化。

\subsection*{一般线性映射的分解}
\addcontentsline{toc}{subsection}{一般线性映射的分解}

极分解与奇异值分解更进一步：它们不要求映射的定义域与陪域相同。

若：

\[
T:V\to W,
\]

其中 $V,W$ 是有限维内积空间，则奇异值分解可以在输入空间和输出空间中分别选取规范正交基，使 $T$ 只沿对应方向进行缩放。

极分解则将 $T$ 写成：

\[
T=U|T|,
\qquad
|T|=(T^*T)^{1/2}.
\]

其中 $|T|$ 是定义在输入空间上的正半定自伴算子；$U$ 一般是部分等距算子。只有在适当的单射、满射或可逆条件下，$U$ 才会成为真正的等距同构。

\section*{本章的结构}
\addcontentsline{toc}{section}{本章的结构}

本章先讨论一般实、复向量空间上的算子：不变子空间、特征结构、上三角化、对角化、Jordan 标准型与复化。

之后再讨论内积空间：伴随、正规算子、自伴算子、谱定理、极分解与奇异值分解。
